% glossary.tex -- Module 7: Glossary (generated by
% Terms merged and de-duplicated from the glossaries of ch1-ch4; sorted A-Z.
\unnumberedchapter[\faBook]{Glossary}\label{bm:glossary}

Every technical word in this book, from A to Z. The meaning is in simple English. The
last column gives the usual term in Chinese physics textbooks. The number after each
term is the chapter where it is explained.

\begin{glossarytable}
  \gterm{Absolute temperature ($T$)\ {\footnotesize\textcolor{hscmuted}{\mdseries(3)}}}{temperature in kelvin, measured from absolute zero: $T(\unit{K}) = T(\unit{\degreeCelsius}) + 273$}{热力学温度（绝对温度）}
  \gterm{Absorption spectrum\ {\footnotesize\textcolor{hscmuted}{\mdseries(1)}}}{a continuous spectrum crossed by dark lines; made when light passes through a cooler, low-density gas}{吸收光谱}
  \gterm{Accelerating charge\ {\footnotesize\textcolor{hscmuted}{\mdseries(1)}}}{a charge that speeds up, slows down or changes direction; it sends out EM waves}{加速运动的电荷}
  \gterm{Accuracy\ {\footnotesize\textcolor{hscmuted}{\mdseries(1)}}}{how close a result is to the true (accepted) value}{准确度}
  \gterm{Aether (luminiferous aether)\ {\footnotesize\textcolor{hscmuted}{\mdseries(1, 4)}}}{an invisible material once thought to fill space and carry light; it does not exist}{（光）以太}
  \gterm{Amplitude\ {\footnotesize\textcolor{hscmuted}{\mdseries(2)}}}{the height of a crest above the middle line of a wave}{振幅}
  \gterm{Ampère--Maxwell law\ {\footnotesize\textcolor{hscmuted}{\mdseries(1)}}}{a magnetic field is made by a current and by a changing electric field}{安培–麦克斯韦定律}
  \gterm{Analyser\ {\footnotesize\textcolor{hscmuted}{\mdseries(2)}}}{a second polarising filter that tests the light from the first one}{检偏器}
  \gterm{Angle of incidence\ {\footnotesize\textcolor{hscmuted}{\mdseries(2)}}}{the angle between the incoming ray and the normal}{入射角}
  \gterm{Angle of reflection\ {\footnotesize\textcolor{hscmuted}{\mdseries(2)}}}{the angle between the reflected ray and the normal}{反射角}
  \gterm{Annihilation\ {\footnotesize\textcolor{hscmuted}{\mdseries(4)}}}{a particle and its antiparticle meet and turn completely into photons}{湮灭}
  \gterm{Anode\ {\footnotesize\textcolor{hscmuted}{\mdseries(3)}}}{the plate in a photocell that collects the electrons}{阳极}
  \gterm{Antenna\ {\footnotesize\textcolor{hscmuted}{\mdseries(1)}}}{a metal rod that sends out or picks up radio waves}{天线}
  \gterm{Antinode\ {\footnotesize\textcolor{hscmuted}{\mdseries(1)}}}{a fixed point in a standing wave where the wave is always strongest}{波腹}
  \gterm{Antiparticle\ {\footnotesize\textcolor{hscmuted}{\mdseries(4)}}}{a particle with the same mass as another particle but the opposite charge}{反粒子}
  \gterm{Atmosphere (of a star)\ {\footnotesize\textcolor{hscmuted}{\mdseries(1)}}}{the cooler, lower-density outer layers of a star, where absorption lines form}{恒星大气}
  \gterm{Atomic clock\ {\footnotesize\textcolor{hscmuted}{\mdseries(1, 4)}}}{a very accurate clock that counts the oscillations of radiation from atoms (for example caesium)}{原子钟}
  \gterm{Black body\ {\footnotesize\textcolor{hscmuted}{\mdseries(1, 3)}}}{an ideal object that absorbs all radiation that hits it and is a perfect emitter; its spectrum depends only on its temperature}{黑体}
  \gterm{Black body radiation\ {\footnotesize\textcolor{hscmuted}{\mdseries(3)}}}{the thermal radiation given off by a black body}{黑体辐射}
  \gterm{Blueshift\ {\footnotesize\textcolor{hscmuted}{\mdseries(1)}}}{spectral lines moved to shorter wavelengths, because the source moves towards us}{蓝移}
  \gterm{Bolometer\ {\footnotesize\textcolor{hscmuted}{\mdseries(3)}}}{a detector whose electrical resistance changes when radiation warms it}{测辐射热计}
  \gterm{Cathode\ {\footnotesize\textcolor{hscmuted}{\mdseries(3)}}}{the metal plate in a photocell that light shines on; it releases electrons}{阴极}
  \gterm{Central maximum\ {\footnotesize\textcolor{hscmuted}{\mdseries(2)}}}{the wide bright band in the middle of a diffraction pattern}{中央亮纹（中央极大）}
  \gterm{Classical physics\ {\footnotesize\textcolor{hscmuted}{\mdseries(3)}}}{physics before about 1900, in which light is a wave and energy can have any value}{经典物理学}
  \gterm{Coherent\ {\footnotesize\textcolor{hscmuted}{\mdseries(2)}}}{waves with the same frequency and a fixed phase difference (they stay in step)}{相干的}
  \gterm{Collimator\ {\footnotesize\textcolor{hscmuted}{\mdseries(1)}}}{the slit-and-lens part of a spectroscope that makes the light rays parallel}{平行光管（准直管）}
  \gterm{Constructive interference\ {\footnotesize\textcolor{hscmuted}{\mdseries(2)}}}{waves meet in phase and add to a bigger wave (bright light)}{相长干涉}
  \gterm{Continuous\ {\footnotesize\textcolor{hscmuted}{\mdseries(3)}}}{able to have any value in a range}{连续的}
  \gterm{Continuous spectrum\ {\footnotesize\textcolor{hscmuted}{\mdseries(1)}}}{a spectrum with all wavelengths and no gaps, from a hot, dense object}{连续光谱（连续谱）}
  \gterm{Corpuscle\ {\footnotesize\textcolor{hscmuted}{\mdseries(2)}}}{Newton's name for a tiny particle of light}{微粒（光微粒）}
  \gterm{Cosmological time dilation\ {\footnotesize\textcolor{hscmuted}{\mdseries(4)}}}{events in very distant galaxies are measured to run slow, by the factor $1 + z$}{宇宙学时间膨胀}
  \gterm{Crest\ {\footnotesize\textcolor{hscmuted}{\mdseries(2)}}}{a highest point of a wave}{波峰}
  \gterm{Decay\ {\footnotesize\textcolor{hscmuted}{\mdseries(4)}}}{an unstable particle breaking up into other particles}{衰变}
  \gterm{Destructive interference\ {\footnotesize\textcolor{hscmuted}{\mdseries(2)}}}{waves meet \ang{180} out of phase and cancel (darkness)}{相消干涉}
  \gterm{Diffraction\ {\footnotesize\textcolor{hscmuted}{\mdseries(2)}}}{the spreading of a wave as it passes through a gap or around an edge}{衍射}
  \gterm{Diffraction grating\ {\footnotesize\textcolor{hscmuted}{\mdseries(1, 2)}}}{a plate with thousands of equally spaced, parallel slits; it spreads light into a spectrum}{衍射光栅}
  \gterm{Direct method\ {\footnotesize\textcolor{hscmuted}{\mdseries(1)}}}{finding $c$ from a measured distance and a measured travel time}{直接测量法}
  \gterm{Discharge tube\ {\footnotesize\textcolor{hscmuted}{\mdseries(1)}}}{a glass tube of low-pressure gas that glows when a high voltage is applied}{气体放电管}
  \gterm{Displacement (of a wave)\ {\footnotesize\textcolor{hscmuted}{\mdseries(2)}}}{how far a wave is above or below its middle line at a point}{位移}
  \gterm{Doppler effect\ {\footnotesize\textcolor{hscmuted}{\mdseries(1)}}}{the change in wavelength when a source and an observer move towards or away from each other}{多普勒效应}
  \gterm{Eclipse (of a moon)\ {\footnotesize\textcolor{hscmuted}{\mdseries(1)}}}{a moon passing into its planet's shadow and disappearing for a while}{卫星食（木卫一食）}
  \gterm{Electric field\ {\footnotesize\textcolor{hscmuted}{\mdseries(1)}}}{a region where a charge feels a force}{电场}
  \gterm{Electromagnetic (EM) spectrum\ {\footnotesize\textcolor{hscmuted}{\mdseries(1)}}}{all EM waves, in order of frequency or wavelength}{电磁波谱}
  \gterm{Electromagnetic (EM) wave\ {\footnotesize\textcolor{hscmuted}{\mdseries(1)}}}{a travelling wave of electric and magnetic fields that change together}{电磁波}
  \gterm{Electromagnetic field\ {\footnotesize\textcolor{hscmuted}{\mdseries(1)}}}{the joined electric and magnetic field; a change in one part makes the other}{电磁场}
  \gterm{Electron volt (eV)\ {\footnotesize\textcolor{hscmuted}{\mdseries(3)}}}{the energy an electron gains moving through \qty{1}{V}: \qty{1}{\eV} $=$ \qty{1.602e-19}{J}}{电子伏特}
  \gterm{Electroscope\ {\footnotesize\textcolor{hscmuted}{\mdseries(3)}}}{a metal rod with a thin gold leaf in a glass case; the leaf stands out when it is charged}{验电器}
  \gterm{Emission spectrum\ {\footnotesize\textcolor{hscmuted}{\mdseries(1)}}}{bright lines on a dark background, from an excited, low-density gas}{发射光谱（明线光谱）}
  \gterm{Energy level\ {\footnotesize\textcolor{hscmuted}{\mdseries(1)}}}{one of the allowed energies of an electron in an atom}{能级}
  \gterm{Envelope\ {\footnotesize\textcolor{hscmuted}{\mdseries(2)}}}{the single-slit curve that dims the outer fringes of a double-slit pattern}{包络线}
  \gterm{Event\ {\footnotesize\textcolor{hscmuted}{\mdseries(4)}}}{something that happens at one place at one instant}{事件}
  \gterm{Excited (atom)\ {\footnotesize\textcolor{hscmuted}{\mdseries(1)}}}{given extra energy, for example by heating or by collisions with fast electrons}{受激的（处于激发态）}
  \gterm{Extrapolate\ {\footnotesize\textcolor{hscmuted}{\mdseries(3)}}}{extend a graph line beyond the measured points}{外推}
  \gterm{Faraday's law\ {\footnotesize\textcolor{hscmuted}{\mdseries(1)}}}{a changing magnetic field makes an electric field (an emf)}{法拉第电磁感应定律}
  \gterm{Frame of reference\ {\footnotesize\textcolor{hscmuted}{\mdseries(4)}}}{a viewpoint, with its own rulers and clocks, for measuring positions, times and motion}{参考系}
  \gterm{Fraunhofer lines\ {\footnotesize\textcolor{hscmuted}{\mdseries(1)}}}{the dark absorption lines in the spectrum of sunlight}{夫琅禾费谱线}
  \gterm{Frequency ($f$)\ {\footnotesize\textcolor{hscmuted}{\mdseries(1, 2)}}}{the number of waves that pass a point each second; unit hertz (Hz)}{频率}
  \gterm{Fringe\ {\footnotesize\textcolor{hscmuted}{\mdseries(2)}}}{one bright or dark band in an interference pattern}{条纹}
  \gterm{Fringe spacing ($\Delta y$)\ {\footnotesize\textcolor{hscmuted}{\mdseries(2)}}}{the distance between the centres of two neighbouring bright fringes}{条纹间距}
  \gterm{Galilean relativity\ {\footnotesize\textcolor{hscmuted}{\mdseries(4)}}}{the laws of motion are the same in every inertial frame}{伽利略相对性原理}
  \gterm{Gauss's law (electric)\ {\footnotesize\textcolor{hscmuted}{\mdseries(1)}}}{charges make electric fields; field lines start on $+$ and end on $-$ charges}{电场的高斯定理}
  \gterm{Gauss's law for magnetism\ {\footnotesize\textcolor{hscmuted}{\mdseries(1)}}}{magnetic field lines always form closed loops; there are no single magnetic poles}{磁场的高斯定理}
  \gterm{Giant star\ {\footnotesize\textcolor{hscmuted}{\mdseries(1)}}}{a huge, swollen star near the end of its life, with low-density outer layers}{巨星}
  \gterm{Gradient\ {\footnotesize\textcolor{hscmuted}{\mdseries(3)}}}{the steepness of a line: rise divided by run}{斜率}
  \gterm{Grating spacing ($d$)\ {\footnotesize\textcolor{hscmuted}{\mdseries(2)}}}{the distance from one slit of a grating to the next}{光栅常量（光栅常数）}
  \gterm{Huygens' principle\ {\footnotesize\textcolor{hscmuted}{\mdseries(2)}}}{every point on a wavefront sends out wavelets; the line touching them is the new wavefront}{惠更斯原理}
  \gterm{In phase\ {\footnotesize\textcolor{hscmuted}{\mdseries(1, 2)}}}{reaching their peaks at the same time and place: crests arrive together with crests}{同相}
  \gterm{Incandescent\ {\footnotesize\textcolor{hscmuted}{\mdseries(1)}}}{glowing because it is hot}{白炽的}
  \gterm{Indirect method\ {\footnotesize\textcolor{hscmuted}{\mdseries(1)}}}{finding $c$ from other measured quantities, such as $f$ and $\lambda$}{间接测量法}
  \gterm{Inertial frame of reference\ {\footnotesize\textcolor{hscmuted}{\mdseries(4)}}}{a frame that is not accelerating: at rest or moving at constant velocity}{惯性参考系}
  \gterm{Intensity ($I$)\ {\footnotesize\textcolor{hscmuted}{\mdseries(2, 3)}}}{the power of light falling on each square metre (\si{W.m^{-2}}); for one colour, it is set by the number of photons each second}{光强（强度）}
  \gterm{Intercept\ {\footnotesize\textcolor{hscmuted}{\mdseries(3)}}}{the value where a line crosses an axis}{截距}
  \gterm{Interference\ {\footnotesize\textcolor{hscmuted}{\mdseries(2)}}}{the pattern made when overlapping waves add or cancel}{干涉}
  \gterm{Interferometer\ {\footnotesize\textcolor{hscmuted}{\mdseries(4)}}}{an instrument that splits and recombines light to compare two light paths}{干涉仪}
  \gterm{Least-bound electron\ {\footnotesize\textcolor{hscmuted}{\mdseries(3)}}}{the electron that is easiest to remove from a metal; it needs exactly the work function}{束缚最弱的电子}
  \gterm{Length contraction\ {\footnotesize\textcolor{hscmuted}{\mdseries(4)}}}{a moving object is measured to be shorter along its direction of motion}{长度收缩（尺缩效应）}
  \gterm{Light clock\ {\footnotesize\textcolor{hscmuted}{\mdseries(4)}}}{a light pulse bouncing between two mirrors; one round trip is one tick}{光钟}
  \gterm{Light-year (ly)\ {\footnotesize\textcolor{hscmuted}{\mdseries(4)}}}{the distance light travels in one year, about \qty{9.46e15}{m}}{光年}
  \gterm{Line broadening\ {\footnotesize\textcolor{hscmuted}{\mdseries(1)}}}{spectral lines becoming wider, from fast rotation or a dense atmosphere}{谱线展宽}
  \gterm{Line of best fit\ {\footnotesize\textcolor{hscmuted}{\mdseries(3)}}}{one straight line drawn as close as possible to all the data points}{最佳拟合直线}
  \gterm{Line of sight\ {\footnotesize\textcolor{hscmuted}{\mdseries(1)}}}{the straight line from an observer to a star; only motion along it gives a Doppler shift}{视线（视线方向）}
  \gterm{Longitudinal wave\ {\footnotesize\textcolor{hscmuted}{\mdseries(2)}}}{a wave that vibrates along its direction of travel, like sound}{纵波}
  \gterm{Lorentz factor ($\gamma$)\ {\footnotesize\textcolor{hscmuted}{\mdseries(4)}}}{$1/\sqrt{1 - v^2/c^2}$: how many times longer a moving clock's tick is}{洛伦兹因子}
  \gterm{Low-density gas\ {\footnotesize\textcolor{hscmuted}{\mdseries(1)}}}{a gas whose atoms are far apart}{稀薄气体（低密度气体）}
  \gterm{Magnetic field\ {\footnotesize\textcolor{hscmuted}{\mdseries(1)}}}{a region where a magnet or a moving charge feels a force}{磁场}
  \gterm{Main-sequence star\ {\footnotesize\textcolor{hscmuted}{\mdseries(1)}}}{an ordinary star, like the Sun, in the long middle part of its life}{主序星}
  \gterm{Malus' law\ {\footnotesize\textcolor{hscmuted}{\mdseries(2)}}}{the intensity through a polariser is $I = I_\mathrm{max}\cos^2\theta$}{马吕斯定律}
  \gterm{Mass defect ($\Delta m$)\ {\footnotesize\textcolor{hscmuted}{\mdseries(4)}}}{mass before minus mass after a process: the mass that becomes energy}{质量亏损}
  \gterm{Mass--energy equivalence\ {\footnotesize\textcolor{hscmuted}{\mdseries(4)}}}{mass and energy are two forms of the same thing: $E = mc^2$}{质能等价（质能方程）}
  \gterm{Maximum kinetic energy ($K_\mathrm{max}$)\ {\footnotesize\textcolor{hscmuted}{\mdseries(3)}}}{the kinetic energy of the fastest photoelectrons}{（光电子的）最大初动能}
  \gterm{Mean lifetime\ {\footnotesize\textcolor{hscmuted}{\mdseries(4)}}}{the average time an unstable particle lives before it decays}{平均寿命}
  \gterm{Medium\ {\footnotesize\textcolor{hscmuted}{\mdseries(4)}}}{a substance that a wave travels through, as sound travels through air}{介质}
  \gterm{Microammeter\ {\footnotesize\textcolor{hscmuted}{\mdseries(3)}}}{a meter for very small currents, in millionths of an ampere}{微安表}
  \gterm{Model\ {\footnotesize\textcolor{hscmuted}{\mdseries(2)}}}{a simple picture that explains what we see and predicts new results}{模型}
  \gterm{Monochromatic\ {\footnotesize\textcolor{hscmuted}{\mdseries(2)}}}{of one colour, so of one wavelength}{单色的}
  \gterm{Muon\ {\footnotesize\textcolor{hscmuted}{\mdseries(4)}}}{an unstable particle like a heavy electron, made by cosmic rays high in the atmosphere}{μ子（缪子）}
  \gterm{Non-inertial frame\ {\footnotesize\textcolor{hscmuted}{\mdseries(4)}}}{a frame that is accelerating: speeding up, slowing down or turning}{非惯性参考系}
  \gterm{Normal\ {\footnotesize\textcolor{hscmuted}{\mdseries(2)}}}{a line at \ang{90} to a surface; ray angles are measured from it}{法线}
  \gterm{Nuclear fusion\ {\footnotesize\textcolor{hscmuted}{\mdseries(4)}}}{light nuclei joining to make a heavier nucleus, releasing energy}{核聚变}
  \gterm{Nucleus (plural: nuclei)\ {\footnotesize\textcolor{hscmuted}{\mdseries(4)}}}{the tiny centre of an atom, made of protons and neutrons}{原子核}
  \gterm{Null result\ {\footnotesize\textcolor{hscmuted}{\mdseries(4)}}}{an experiment that does not find the effect it was looking for}{零结果}
  \gterm{Observer\ {\footnotesize\textcolor{hscmuted}{\mdseries(4)}}}{a person or instrument that makes measurements in one frame of reference}{观察者}
  \gterm{Order ($m$)\ {\footnotesize\textcolor{hscmuted}{\mdseries(2)}}}{the number of a fringe, counted out from the centre (bright: $m = 0, 1, 2, \ldots$; dark: $m = 1, 2, \ldots$)}{级次}
  \gterm{Oscillator\ {\footnotesize\textcolor{hscmuted}{\mdseries(3)}}}{a vibrating charged particle in the walls of a black body (Planck's model)}{振子}
  \gterm{Out of phase (\ang{180})\ {\footnotesize\textcolor{hscmuted}{\mdseries(2)}}}{crests of one wave arrive with troughs of the other}{反相}
  \gterm{Particle accelerator\ {\footnotesize\textcolor{hscmuted}{\mdseries(4)}}}{a machine that speeds up charged particles with electric fields and steers them with magnetic fields}{粒子加速器}
  \gterm{Path difference ($\delta$)\ {\footnotesize\textcolor{hscmuted}{\mdseries(2)}}}{how much further one wave travels than the other}{波程差（路程差）}
  \gterm{Peak wavelength ($\lambda_\mathrm{max}$)\ {\footnotesize\textcolor{hscmuted}{\mdseries(1, 3)}}}{the wavelength where a black body curve is highest}{峰值波长}
  \gterm{Permeability of free space ($\mu_0$)\ {\footnotesize\textcolor{hscmuted}{\mdseries(1)}}}{a constant that sets how strong magnetic fields are in a vacuum}{真空磁导率}
  \gterm{Permittivity of free space ($\varepsilon_0$)\ {\footnotesize\textcolor{hscmuted}{\mdseries(1)}}}{a constant that sets how strong electric fields are in a vacuum}{真空介电常量（真空电容率）}
  \gterm{Phase\ {\footnotesize\textcolor{hscmuted}{\mdseries(2)}}}{where a wave is in its cycle: at a crest, at a trough, or in between}{相位}
  \gterm{Photocell\ {\footnotesize\textcolor{hscmuted}{\mdseries(3)}}}{the vacuum tube, with a metal cathode and an anode, used in photoelectric experiments}{光电管}
  \gterm{Photocurrent\ {\footnotesize\textcolor{hscmuted}{\mdseries(3)}}}{the current made by photoelectrons}{光电流}
  \gterm{Photoelectric effect\ {\footnotesize\textcolor{hscmuted}{\mdseries(3)}}}{electrons are released from a metal surface when light shines on it}{光电效应}
  \gterm{Photoelectron\ {\footnotesize\textcolor{hscmuted}{\mdseries(3)}}}{an electron released from a metal by light}{光电子}
  \gterm{Photon\ {\footnotesize\textcolor{hscmuted}{\mdseries(1, 3)}}}{one packet (quantum) of light energy, $E = hf$}{光子}
  \gterm{Planck's constant ($h$)\ {\footnotesize\textcolor{hscmuted}{\mdseries(3)}}}{\qty{6.626e-34}{J.s}; it links the energy of a quantum to its frequency}{普朗克常量}
  \gterm{Plane-polarised light\ {\footnotesize\textcolor{hscmuted}{\mdseries(2)}}}{light whose electric field vibrates in one direction only}{线偏振光（平面偏振光）}
  \gterm{Polarisation (polarised)\ {\footnotesize\textcolor{hscmuted}{\mdseries(2)}}}{a transverse wave is polarised when it vibrates in one direction (one plane) only}{偏振}
  \gterm{Polariser\ {\footnotesize\textcolor{hscmuted}{\mdseries(2)}}}{a filter that passes only the vibration along its transmission axis}{起偏器（偏振片）}
  \gterm{Positron ($\mathrm{e}^+$)\ {\footnotesize\textcolor{hscmuted}{\mdseries(4)}}}{the antiparticle of the electron: same mass, positive charge}{正电子}
  \gterm{Postulate\ {\footnotesize\textcolor{hscmuted}{\mdseries(4)}}}{a starting statement that a theory takes as true}{基本假设（公设）}
  \gterm{Precision\ {\footnotesize\textcolor{hscmuted}{\mdseries(1)}}}{how close repeated results are to each other}{精密度}
  \gterm{Pressure broadening\ {\footnotesize\textcolor{hscmuted}{\mdseries(1)}}}{wider lines, because atoms in a dense gas collide often and their energy levels are disturbed}{压力展宽（碰撞展宽）}
  \gterm{Proper length ($l_0$)\ {\footnotesize\textcolor{hscmuted}{\mdseries(4)}}}{the length of an object measured in the frame where it is at rest}{固有长度（原长）}
  \gterm{Proper time ($t_0$)\ {\footnotesize\textcolor{hscmuted}{\mdseries(4)}}}{the time between two events, measured by a clock that is present at both events}{固有时（原时）}
  \gterm{Proton--proton chain\ {\footnotesize\textcolor{hscmuted}{\mdseries(4)}}}{the chain of fusion reactions in which the Sun turns four hydrogen nuclei into one helium-4 nucleus}{质子–质子链}
  \gterm{Quantised (discrete)\ {\footnotesize\textcolor{hscmuted}{\mdseries(3)}}}{can only have certain separate values, in steps}{量子化的（分立的）}
  \gterm{Quantum (plural: quanta)\ {\footnotesize\textcolor{hscmuted}{\mdseries(3)}}}{one packet of energy, of size $E = hf$}{量子}
  \gterm{Random error\ {\footnotesize\textcolor{hscmuted}{\mdseries(1)}}}{small, unpredictable changes between repeated readings; averaging reduces it}{随机误差（偶然误差）}
  \gterm{Rayleigh--Jeans law\ {\footnotesize\textcolor{hscmuted}{\mdseries(3)}}}{the classical prediction for the black body curve; correct only at long wavelengths}{瑞利–金斯公式}
  \gterm{Redshift\ {\footnotesize\textcolor{hscmuted}{\mdseries(1, 4)}}}{spectral lines moved to longer wavelengths, for example because the source moves away from us}{红移}
  \gterm{Refraction\ {\footnotesize\textcolor{hscmuted}{\mdseries(2)}}}{light changing direction as it goes from one material into another}{折射}
  \gterm{Refractive index ($n$)\ {\footnotesize\textcolor{hscmuted}{\mdseries(2)}}}{how many times slower light is in a material than in a vacuum}{折射率}
  \gterm{Relativistic momentum ($p_\mathrm{v}$)\ {\footnotesize\textcolor{hscmuted}{\mdseries(4)}}}{the momentum that is correct at any speed (on the formula sheet); it grows without limit as $v$ approaches $c$}{相对论动量}
  \gterm{Relativity of simultaneity\ {\footnotesize\textcolor{hscmuted}{\mdseries(4)}}}{events at different places that are simultaneous in one inertial frame are not simultaneous in another}{同时的相对性}
  \gterm{Reliability\ {\footnotesize\textcolor{hscmuted}{\mdseries(1)}}}{repeating the experiment gives consistent results}{可靠性}
  \gterm{Repulsive force\ {\footnotesize\textcolor{hscmuted}{\mdseries(2)}}}{a push away from a surface (Newton's explanation of reflection)}{斥力}
  \gterm{Rest energy\ {\footnotesize\textcolor{hscmuted}{\mdseries(4)}}}{the energy an object has because of its mass, $E = mc^2$}{静能（静止能量）}
  \gterm{Rest mass ($m_0$)\ {\footnotesize\textcolor{hscmuted}{\mdseries(4)}}}{the mass measured by an observer for whom the object is at rest}{静止质量}
  \gterm{Rotational velocity (of a star)\ {\footnotesize\textcolor{hscmuted}{\mdseries(1)}}}{how fast a star spins: the speed of its surface as it turns}{自转速度}
  \gterm{Saturation current\ {\footnotesize\textcolor{hscmuted}{\mdseries(3)}}}{the largest photocurrent, when every released electron reaches the anode}{饱和电流（饱和光电流）}
  \gterm{Simultaneous\ {\footnotesize\textcolor{hscmuted}{\mdseries(4)}}}{happening at the same time}{同时（的）}
  \gterm{Slit separation ($d$)\ {\footnotesize\textcolor{hscmuted}{\mdseries(2)}}}{the distance between the centres of two slits}{双缝间距（缝间距）}
  \gterm{Slit width ($a$)\ {\footnotesize\textcolor{hscmuted}{\mdseries(2)}}}{how wide one slit is}{缝宽}
  \gterm{Spark gap\ {\footnotesize\textcolor{hscmuted}{\mdseries(3)}}}{a small air gap between two metal balls that a spark jumps across}{火花隙}
  \gterm{Special relativity\ {\footnotesize\textcolor{hscmuted}{\mdseries(4)}}}{Einstein's theory of how space and time behave for observers in inertial frames}{狭义相对论}
  \gterm{Spectral line\ {\footnotesize\textcolor{hscmuted}{\mdseries(1)}}}{one wavelength that is much brighter or darker than its neighbours}{谱线}
  \gterm{Spectroscope\ {\footnotesize\textcolor{hscmuted}{\mdseries(1, 2, 4)}}}{an instrument that spreads light into a spectrum so you can study it}{分光镜（光谱仪）}
  \gterm{Spectroscopy\ {\footnotesize\textcolor{hscmuted}{\mdseries(1)}}}{using spectra to learn about a light source}{光谱学（光谱分析）}
  \gterm{Spectrum (plural: spectra)\ {\footnotesize\textcolor{hscmuted}{\mdseries(1, 2)}}}{light spread out in order of wavelength}{光谱}
  \gterm{Speed of light ($c$)\ {\footnotesize\textcolor{hscmuted}{\mdseries(1)}}}{the speed of all EM waves in a vacuum, exactly \qty{299792458}{m.s^{-1}}}{光速}
  \gterm{Standing wave\ {\footnotesize\textcolor{hscmuted}{\mdseries(1)}}}{a wave pattern that stays in one place instead of travelling}{驻波}
  \gterm{Stopping voltage ($V_0$)\ {\footnotesize\textcolor{hscmuted}{\mdseries(3)}}}{the smallest reverse voltage that stops every photoelectron, so the current is zero}{遏止电压}
  \gterm{Superposition\ {\footnotesize\textcolor{hscmuted}{\mdseries(2)}}}{when two waves meet, their displacements add}{叠加（波的叠加）}
  \gterm{Synchrotron\ {\footnotesize\textcolor{hscmuted}{\mdseries(4)}}}{a ring-shaped accelerator whose magnets get stronger as the particles gain momentum}{同步加速器}
  \gterm{Thermal radiation\ {\footnotesize\textcolor{hscmuted}{\mdseries(3)}}}{EM radiation that an object gives off because of its temperature}{热辐射}
  \gterm{Thought experiment\ {\footnotesize\textcolor{hscmuted}{\mdseries(4)}}}{an imagined experiment, worked through with logic, to see what a theory predicts}{思想实验}
  \gterm{Threshold frequency ($f_0$)\ {\footnotesize\textcolor{hscmuted}{\mdseries(3)}}}{the lowest frequency of light that releases electrons from a metal}{截止频率（极限频率）}
  \gterm{Threshold wavelength ($\lambda_0$)\ {\footnotesize\textcolor{hscmuted}{\mdseries(3)}}}{the longest wavelength of light that releases electrons from a metal}{截止波长（极限波长）}
  \gterm{Time dilation\ {\footnotesize\textcolor{hscmuted}{\mdseries(4)}}}{a moving clock is measured to run slow}{时间膨胀（时间延缓）}
  \gterm{Toothed wheel\ {\footnotesize\textcolor{hscmuted}{\mdseries(1)}}}{a wheel with equal teeth and gaps around its edge, used by Fizeau to time light}{齿轮（菲索齿轮）}
  \gterm{Translational velocity (of a star)\ {\footnotesize\textcolor{hscmuted}{\mdseries(1)}}}{how fast the whole star moves through space}{平动速度}
  \gterm{Transmission axis\ {\footnotesize\textcolor{hscmuted}{\mdseries(2)}}}{the direction of vibration that a polarising filter lets through}{透振方向（透光轴）}
  \gterm{Transverse wave\ {\footnotesize\textcolor{hscmuted}{\mdseries(1, 2)}}}{a wave whose vibrations are at \ang{90} to its direction of travel}{横波}
  \gterm{Trough\ {\footnotesize\textcolor{hscmuted}{\mdseries(2)}}}{a lowest point of a wave}{波谷}
  \gterm{Type Ia supernova\ {\footnotesize\textcolor{hscmuted}{\mdseries(4)}}}{an exploding star whose brightness rises and falls in nearly the same way every time}{Ia型超新星}
  \gterm{Ultraviolet catastrophe\ {\footnotesize\textcolor{hscmuted}{\mdseries(3)}}}{the classical prediction of infinite intensity at short wavelengths, which does not happen}{紫外灾难}
  \gterm{Unification\ {\footnotesize\textcolor{hscmuted}{\mdseries(1)}}}{showing that two things that look different are parts of one thing}{统一}
  \gterm{Unpolarised light\ {\footnotesize\textcolor{hscmuted}{\mdseries(2)}}}{light that vibrates in all directions across the beam, changing at random}{非偏振光（自然光）}
  \gterm{Unstable\ {\footnotesize\textcolor{hscmuted}{\mdseries(4)}}}{breaks up into other particles after a short time}{不稳定的}
  \gterm{Validity\ {\footnotesize\textcolor{hscmuted}{\mdseries(1)}}}{the method really tests what it claims to, with other variables controlled}{有效性}
  \gterm{Wavefront\ {\footnotesize\textcolor{hscmuted}{\mdseries(2)}}}{a line joining points of a wave that are in step, such as all the crests}{波阵面（波面）}
  \gterm{Wavelength ($\lambda$)\ {\footnotesize\textcolor{hscmuted}{\mdseries(1, 2)}}}{the length of one complete wave, for example from crest to crest; unit metre (m)}{波长}
  \gterm{Wavelet\ {\footnotesize\textcolor{hscmuted}{\mdseries(2)}}}{a small circular wave sent out from one point on a wavefront}{子波}
  \gterm{Wave--particle duality\ {\footnotesize\textcolor{hscmuted}{\mdseries(3)}}}{light acts like a wave in some experiments and like particles in others}{波粒二象性}
  \gterm{Wien's displacement constant ($b$)\ {\footnotesize\textcolor{hscmuted}{\mdseries(3)}}}{\qty{2.898e-3}{m.K}, the constant in Wien's law}{维恩位移常量}
  \gterm{Wien's displacement law\ {\footnotesize\textcolor{hscmuted}{\mdseries(1, 3)}}}{the peak wavelength of a black body is inversely proportional to its absolute temperature: $\lambda_\mathrm{max} = \dfrac{b}{T}$}{维恩位移定律}
  \gterm{Work function ($\phi$)\ {\footnotesize\textcolor{hscmuted}{\mdseries(3)}}}{the smallest energy needed to remove an electron from a metal's surface}{逸出功}
  \gterm{Young's double-slit experiment\ {\footnotesize\textcolor{hscmuted}{\mdseries(2)}}}{light through two close slits makes bright and dark fringes, showing interference}{杨氏双缝实验}
\end{glossarytable}
